\section{The actual increment, not another Forge architecture}
\label{sec:delta}

\subsection{What the reviewed baseline already contains}
The baseline is Forge commit
\code{674521027d968d59f7b83220ed52304a85cb55e2}, particularly the merged article's
structural-search section, provenance appendix, status ledger, and review of
Leant and Djex. It already covers induction planning, accumulator
generalization, demand-directed theorem use, polynomial recurrence synthesis,
algebraic certificates, and a carefully delimited acceptance boundary
\cite{forge,forge-structural,forge-status,forge-neighbors}. These are dependencies
of this proposal, not its contributions.

The relevant unfinished step is narrower and more consequential. The implemented
state-machine invariant certificate described in the merged article requires
$I(F(x))=I(x)$. Its prospective generalization allows multiple invariants and
polynomial multipliers, but correctly warns that discovering both at once is a
bilinear problem. The first contribution below supplies a different search
organization: demand a \emph{closed space of observations}, discover it by exact
linear dependence, and then prove that its generators all vanish on the
reachable states. No simultaneous unknown invariant/multiplier system is solved.

The second relevant gap is already named in the repository: negative results
from an intuitionistic inhabitation decision procedure lack a separately
checkable certificate. The proposal here does not claim that identifying this
gap is new. It provides an implemented \emph{semantic alternative} to logging an
entire unsuccessful LJT search: a finite Kripke structure and a small forcing
checker.

\begin{table}[ht]
\centering\small
\begin{tabular}{@{}P{0.27\textwidth}P{0.30\textwidth}P{0.35\textwidth}@{}}
\toprule
Existing ingredient & Increment in this package & Resulting capability \\
\midrule
Conserved polynomial invariant & Control-indexed coefficient-pullback closure & Mutually supporting relational invariants, discovered rather than supplied \\
Polynomial recurrence goals & Fresh symbolic payloads at every edge & All finite words over an infinite rational input alphabet \\
Typed candidate plus bounded behavioral tests & Product-machine proof obligation & Universal fold equivalence in an explicit fragment \\
Suggested failed-search certificate & Finite Kripke model and independent forcing & Replayable IPC nonderivability diagnostics \\
Uncompiled certificate sources & 222 generated local identity obligations & A concrete next compiler gate, not a claim of Lean acceptance \\
\bottomrule
\end{tabular}
\caption{The delta relative to the reviewed merged design. The last row is an integration artifact, not a new logical method.}
\end{table}

\subsection{Two recommendations}
\textbf{Implement the observable-closure worker first.} It addresses a positive
proof capability directly: statements about arbitrary finite executions,
streaming algorithms, accumulator programs, and equivalence of two folds. Its
certificate needs only exact polynomial identities plus induction, and it can
reuse existing Lean arithmetic reconstruction.

\textbf{Implement the Kripke worker as a diagnostic lane.} It is valuable when a
user asks whether constructive propositional reasoning is sufficient. It must
not close the original goal with its negation. Its strongest appropriate output
is a checked theorem that the corresponding \emph{object-logic sequent} has no
derivation.

These are not substitutes for each other. A rational execution counterexample
can refute a particular universally quantified machine property. An IPC
countermodel instead classifies a restricted proof language. An interface that
uses one generic ``refuted'' label for both would erase the most important
semantic distinction in the proposal.

\subsection{Relation to established work}
The algebraic worker belongs to the tradition of exact invariant analysis,
including Karr's affine relations, and to the linear-algebraic treatment of
observables under composition. Finite invariant spaces of Koopman observables
are an established perspective on nonlinear dynamics; this proposal uses exact
polynomial pullbacks for discrete, symbolic-input proof obligations, not
numerically learned approximate observables \cite{karr,koopman}. Register
automata with linear arithmetic supply closely related algorithmic context
\cite{register}. Neither invariant subspaces nor rational-machine verification
is claimed to originate here.

Kripke semantics is likewise established. Lean formalizations of intuitionistic
propositional logic already provide soundness and completeness infrastructure
\cite{trufas}. Mathlib's \code{itauto} already implements a complete positive
intuitionistic propositional procedure for its supported syntax
\cite{itauto}. The addition is a reusable negative \emph{certificate boundary},
not a new claim to solve positive propositional tautologies.

\subsection{Evidence labels used throughout}
\emph{Proved mathematically} means a proof is given in this article.
\emph{Executed} means the supplied code actually ran and has retained results.
\emph{Generated, uncompiled} means Lean text was emitted but no Lean compiler
checked it. \emph{Proposed} means an implementation plan, not completed code.
These labels must not be collapsed into a single word such as ``verified.''
The supplied Python checker is an executable implementation of the mathematics;
it is not itself formally verified.
